This paper examines how tighter enforcement of cropland protection shapes local land-use behavior in a centralized administrative system characterized by multitask incentives and imperfect monitoring. Leveraging cross-county differences in exposure to cropland protection pressure, the analysis shows that the post-2017 policy tightening generates responses along both the extensive and intensive margins of land use. Counties subject to higher pressure are more likely to expand cropland area, yet the land brought into cultivation is of lower soil suitability and is cultivated less intensively, with elevated fallow rates that persist over time. These patterns suggest that stricter enforcement of observable quantity targets can induce distortions in land allocation and land use when monitoring of land quality and utilization remains limited.

The consequences of these distortions extend beyond short-run mismatches between land endowments and land-use categories. The reclamation of marginal or low-quality land weakens agricultural productivity not only through inferior land characteristics but also through reduced cultivation intensity, higher abandonment, and inefficient input use. Over time, these channels translate into economic losses that are not reflected in aggregate cropland area statistics. As a result, policies that succeed in stabilizing cropland quantity may nevertheless fail to achieve broader productivity objectives if they redirect land use toward plots with limited productive potential.

The findings also inform the design of policy implementation. A common response to execution distortions is to strengthen monitoring to additional dimensions, such as land quality and utilization, yet this approach may be constrained by administrative capacity. This study instead highlights the role of institutional arrangements, particularly the restricted functioning of cropland quota markets. Because cropland replenishment indicators do not circulate freely across prefectures, local governments facing binding targets have limited adjustment margins beyond local land reclamation. Greater flexibility and integration in cropland indicator markets would enlarge the set of feasible compliance options, reducing the likelihood that multitask pressure results in distortive implementation. More broadly, the results underscore that mitigating misallocation under multitask governance requires not only tighter oversight, but also institutional designs that align incentives with policy objectives under realistic monitoring constraints.
